\subsection{Preliminary Consistency Evaluation}
\label{sec:consistency_test}

Building on the consistency framework described above, we conducted a preliminary evaluation to test whether structurally consistent law modifications improve scientific law discovery relative to the original NewtonBench setup. Concretely, we compared \texttt{consistent} and \texttt{inconsistent} benchmark variants under otherwise matched experimental conditions, using the original prompt setting as a controlled baseline. The main trend is encouraging but not uniform. For the metaphysically shifted tasks (\texttt{v0}--\texttt{v2}), consistency yielded a clear improvement for the code-assisted agent: with \texttt{gpt41mini}, mean exact accuracy increased from 11.4\% to 18.4\% across 324 tasks. In contrast, the same intervention produced only a marginal gain for the vanilla agent (from 4.1\% to 4.9\%). At the same time, the effect did not generalize to the \texttt{v\_unchanged} control condition, where consistency provided no benefit and slightly reduced performance in the code-assisted setting. This pattern suggests that consistency is not a generic performance booster; rather, it appears most useful when the benchmark laws themselves instantiate coordinated structural shifts.

We therefore interpret these results as preliminary evidence that consistency-preserving transformations can make the benchmark more legible to stronger agents, especially when those agents are equipped to externalize intermediate reasoning. At the same time, the observed gains remain conditional: they are concentrated in the altered-law regimes and do not imply universal improvement across all task types. These findings motivated our subsequent experimental step. To make the intended cross-law structure more salient, we next modified the prompt format so that the agent is explicitly reminded of what has already been discovered about the simulated universe. The design and analysis of that prompt-level intervention are deferred to the following subsection.
